\section{Diseño del Subsistema de Gestión}

El subsistema de gestión recibe las imágenes capturadas en los accesos, coordina
su procesamiento y convierte el resultado en información persistente,
consultable y auditable. Su diseño se rige por el principio de responsabilidad
única: la base de datos es manipulada exclusivamente por los servicios, y los
servicios no conocen los detalles internos de los modelos de visión. De este
modo, el reentrenamiento o la sustitución de un modelo no impacta sobre las
reglas de negocio ni sobre la estructura de los datos.

\subsection{Arquitectura del subsistema de gestión}

El subsistema se organiza en cuatro componentes. La interfaz de programación de
aplicaciones (API REST) expone los servicios del sistema, valida la estructura
de las solicitudes y constituye el único punto de acceso a los datos. El
notificador mantiene los canales persistentes con los clientes web y propaga las
actualizaciones de estado. La cola de tareas desacopla la recepción de una
solicitud de su procesamiento, y el trabajador de inferencia consume las tareas
encoladas y deposita sus resultados en la capa de persistencia. Esta separación
permite que las reglas de negocio evolucionen sin alterar el motor de
inferencia, y que el motor de inferencia se optimice sin reescribir la lógica de
gestión.

\subsection{Modelo entidad-relación}

El modelo de datos se estructura en seis entidades. La entidad
\textit{vehículo} representa a cada automóvil observado en los accesos y
conserva los atributos visuales reconocidos en su primera aparición. La entidad
\textit{sesión} representa una visita al recinto: se crea en el cruce de ingreso
y se cierra en el cruce de salida, registrando la duración de la permanencia, el
estado del cierre y, cuando el cierre es manual, el operador y el motivo. La
entidad \textit{detección} constituye la bitácora de inferencia: conserva cada
observación producida por el modelo, con su clase, confianza, latencia y
trayectoria asociada. A estas tres se suman la entidad \textit{cámara}, que
registra los puntos de acceso configurados con su puerta, su URL de captura y su
estado de habilitación; la entidad \textit{usuario}, que almacena la credencial,
el rol y el estado de la cuenta; y la entidad \textit{vehículo autorizado}, que
materializa el padrón institucional contra el cual se contrasta cada matrícula
reconocida para el control de autorización (RF-10). La Figura
\ref{fig:modelo_er} presenta la totalidad de las entidades y sus relaciones con
la cardinalidad correspondiente. Las relaciones trazadas con línea continua
corresponden a claves foráneas declaradas en el esquema; las trazadas con línea
discontinua son asociaciones lógicas por coincidencia de identificador o de
matrícula, que el sistema resuelve en la capa de aplicación.

\begin{figure}[H]
    \centering
    \begin{tikzpicture}
        % *** Entidades (nombre conceptual y nombre físico de la tabla) ***
        \spEntityTable{entVehiculo}{-6.0,2.2}{4.0cm}{purpura}
            {vehículo}{vehicles}
            {id (PK)\\license\_plate (ix)\\license\_plate\_confidence\\
             vehicle\_type, color\\+ su confianza\\entry\_image\_path\\
             created\_at, updated\_at}

        \spEntityTable{entSesion}{-0.3,2.2}{4.4cm}{purpura}
            {sesión}{parking\_sessions}
            {id (PK)\\vehicle\_id (FK, ix)\\entry\_gate, exit\_gate\\
             entry\_time, exit\_time\\duration\_seconds\\
             status (ix)\\close\_reason, closed\_by}

        \spEntityTable{entDeteccion}{5.6,2.2}{4.8cm}{matblue}
            {detección}{detection\_events}
            {id (PK)\\session\_id (FK, ix)\\camera\_id (ix), frame\_timestamp\\
             model\_name, model\_version\\object\_class, confidence\\
             bbox (JSONB), crop\_image\_path\\track\_id, inference\_latency\_ms}

        \spEntityTable{entCamara}{5.6,-2.9}{4.8cm}{cielo}
            {cámara}{cameras}
            {id (PK)\\camera\_id (único, ix)\\gate, label, rtsp\_url\\
             is\_active, created\_at}

        \spEntityTable{entAutorizado}{-6.0,-2.9}{4.0cm}{rosa}
            {vehículo autorizado}{authorized\_vehicles}
            {id (PK)\\license\_plate (único, ix)\\owner\_name\\vehicle\_type, color\\
             notes, is\_active\\created\_at, updated\_at}

        \spEntityTable{entUsuario}{-0.3,-2.9}{4.4cm}{rosa}
            {usuario}{users}
            {id (PK)\\username (único, ix)\\password\_hash\\
             role (admin, operator, owner)\\is\_active\\
             created\_at, last\_login\_at}

        % *** Relaciones por clave foránea ***
        \draw[spflow] (entVehiculo.east) --
            node[above, font=\tiny\bfseries, text=black!70] {1:N}
            (entSesion.west);

        \draw[spflow] (entSesion.east) --
            node[above, font=\tiny\bfseries, text=black!70] {1:N}
            (entDeteccion.west);

        % *** Relaciones lógicas ***
        \draw[spflow, dashed] (entAutorizado.north) --
            node[right, font=\tiny\bfseries, text=black!60] {1:N (por matrícula)}
            (entVehiculo.south);

        \draw[spflow, dashed] (entUsuario.north) --
            node[right, font=\tiny\bfseries, text=black!60] {0..1:N (cierre manual)}
            (entSesion.south);

        \draw[spflow, dashed] (entCamara.north) --
            node[right, font=\tiny\bfseries, text=black!60] {1:N (por identificador)}
            (entDeteccion.south);

    \end{tikzpicture}
    \caption{Modelo entidad-relación del subsistema de gestión. Cada entidad
    indica su nombre conceptual y, en tipografía monoespaciada, el nombre físico
    de la tabla en la base de datos.}
    \label{fig:modelo_er}
\end{figure}

La Tabla \ref{tab:diccionario_datos} detalla el diccionario de datos. Cada bloque
identifica la entidad con su nombre conceptual y el nombre físico de la tabla, y
cada fila emplea el nombre real de la columna tal como se declara en las
migraciones, de modo que el esquema pueda verificarse contra el código sin
traducciones intermedias. Las claves foráneas preservan la consistencia
referencial y permiten reconstruir la trazabilidad completa de una visita, desde
el vehículo identificado hasta cada una de las observaciones que sustentaron su
identificación. El esquema se gestiona mediante migraciones versionadas, de modo
que toda modificación estructural quede registrada y pueda reproducirse sobre
cualquier entorno de despliegue.

\begin{longtable}{
        Q{0.28}
        Q{0.16}
        P{0.38}
    }
    \caption{Diccionario de datos del subsistema de gestión}
    \label{tab:diccionario_datos}\\

    \toprule
    \rowcolor{gray!20}
    \textbf{Columna} &
    \textbf{Tipo} &
    \textbf{Restricción} \\
    \midrule
    \endfirsthead

    \multicolumn{3}{c}{\textit{Continuación de la Tabla
    \ref{tab:diccionario_datos}}} \\
    \toprule
    \rowcolor{gray!20}
    \textbf{Columna} &
    \textbf{Tipo} &
    \textbf{Restricción} \\
    \midrule
    \endhead

    \multicolumn{3}{l}{\normalfont\textbf{vehículo} --- \texttt{vehicles}} \\
    \midrule
    \texttt{id} & UUID & Clave primaria. \\
    \texttt{license\_plate} & VARCHAR(20) & Indexado; nulo si no pudo reconocerse. \\
    \texttt{license\_plate\_confidence} & REAL & Puede ser nulo. \\
    \texttt{vehicle\_type} & VARCHAR(50) & Puede ser nulo. \\
    \texttt{vehicle\_type\_confidence} & REAL & Puede ser nulo. \\
    \texttt{color} & VARCHAR(30) & Puede ser nulo. \\
    \texttt{color\_confidence} & REAL & Puede ser nulo. \\
    \texttt{entry\_image\_path} & TEXT & Puede ser nulo. \\
    \texttt{created\_at} & TIMESTAMP & No nulo. \\
    \texttt{updated\_at} & TIMESTAMP & No nulo; se actualiza en cada modificación. \\
    \midrule

    \multicolumn{3}{l}{\normalfont\textbf{sesión} --- \texttt{parking\_sessions}} \\
    \midrule
    \texttt{id} & UUID & Clave primaria. \\
    \texttt{vehicle\_id} & UUID & Clave foránea a \texttt{vehicles}; no nulo e indexado. \\
    \texttt{entry\_gate} & VARCHAR(20) & Puede ser nulo. \\
    \texttt{exit\_gate} & VARCHAR(20) & Puede ser nulo. \\
    \texttt{entry\_time} & TIMESTAMP & No nulo. \\
    \texttt{exit\_time} & TIMESTAMP & Nulo mientras la sesión permanece activa. \\
    \texttt{duration\_seconds} & INTEGER & Segundos de permanencia; calculado al cierre. \\
    \texttt{status} & ENUM & \textit{active}, \textit{completed} o \textit{cancelled};
        indexado. \\
    \texttt{close\_reason} & VARCHAR(50) & Nulo cuando la sesión sigue activa; motivo del cierre. \\
    \texttt{closed\_by} & VARCHAR(50) & Nulo cuando el cierre es automático; operador que cerró. \\
    \midrule

    \multicolumn{3}{l}{\normalfont\textbf{detección} --- \texttt{detection\_events}} \\
    \midrule
    \texttt{id} & UUID & Clave primaria. \\
    \texttt{session\_id} & UUID & Clave foránea a \texttt{parking\_sessions}; puede ser
        nulo e indexado. \\
    \texttt{camera\_id} & VARCHAR(50) & No nulo e indexado. \\
    \texttt{frame\_timestamp} & TIMESTAMP & No nulo. \\
    \texttt{model\_name} & VARCHAR(20) & Identificador del modelo empleado. \\
    \texttt{model\_version} & VARCHAR(20) & Puede ser nulo. \\
    \texttt{object\_class} & VARCHAR(50) & Clase detectada. \\
    \texttt{confidence} & REAL & No nulo. \\
    \texttt{bbox} & JSONB & Coordenadas normalizadas de la caja delimitadora. \\
    \texttt{crop\_image\_path} & TEXT & Puede ser nulo. \\
    \texttt{track\_id} & INTEGER & Identificador de trayectoria asignado por el seguimiento. \\
    \texttt{inference\_latency\_ms} & REAL & Tiempo de inferencia en milisegundos. \\
    \midrule

    \multicolumn{3}{l}{\normalfont\textbf{cámara} --- \texttt{cameras}} \\
    \midrule
    \texttt{id} & UUID & Clave primaria. \\
    \texttt{camera\_id} & VARCHAR(64) & Único, indexado y no nulo; debe coincidir con el
        empleado por el capturador y el trabajador. \\
    \texttt{gate} & VARCHAR(20) & No nulo; punto de acceso asociado. \\
    \texttt{label} & VARCHAR(120) & Puede ser nulo; denominación legible. \\
    \texttt{rtsp\_url} & TEXT & Puede ser nulo; cadena de captura del flujo. \\
    \texttt{is\_active} & BOOLEAN & No nulo; indica si la cámara se espera en operación. \\
    \texttt{created\_at} & TIMESTAMP & No nulo. \\
    \midrule

    \multicolumn{3}{l}{\normalfont\textbf{usuario} --- \texttt{users}} \\
    \midrule
    \texttt{id} & UUID & Clave primaria. \\
    \texttt{username} & VARCHAR(64) & Único, indexado y no nulo. \\
    \texttt{password\_hash} & VARCHAR(255) & No nulo; se almacena únicamente el resumen
        criptográfico, nunca la contraseña. \\
    \texttt{role} & VARCHAR(20) & Restringido por \texttt{CHECK} a \textit{admin},
        \textit{operator} u \textit{owner}. \\
    \texttt{is\_active} & BOOLEAN & No nulo; habilita o suspende la cuenta. \\
    \texttt{created\_at} & TIMESTAMP & No nulo. \\
    \texttt{last\_login\_at} & TIMESTAMP & Puede ser nulo. \\
    \midrule

    \multicolumn{3}{l}{\normalfont\textbf{vehículo autorizado} --- \texttt{authorized\_vehicles}} \\
    \midrule
    \texttt{id} & UUID & Clave primaria. \\
    \texttt{license\_plate} & VARCHAR(20) & Única, indexada y no nulo; sostiene el
        contraste de RF-10. \\
    \texttt{owner\_name} & VARCHAR(120) & Puede ser nulo. \\
    \texttt{vehicle\_type} & VARCHAR(50) & Puede ser nulo. \\
    \texttt{color} & VARCHAR(30) & Puede ser nulo. \\
    \texttt{notes} & TEXT & Puede ser nulo. \\
    \texttt{is\_active} & BOOLEAN & No nulo; una baja lógica conserva el historial de eventos. \\
    \texttt{created\_at} & TIMESTAMP & No nulo. \\
    \texttt{updated\_at} & TIMESTAMP & No nulo; se actualiza en cada modificación. \\*
    \multicolumn{3}{r}{\footnotesize\textit{Fuente: Elaboración propia}} \\*
    \bottomrule
\end{longtable}

Para sostener el desempeño de las consultas sobre un volumen creciente de
registros, el esquema define índices sobre \texttt{license\_plate}, sobre
\texttt{vehicle\_id}, sobre \texttt{status}, sobre \texttt{camera\_id} y sobre
\texttt{session\_id}. A ellos se suman los índices únicos sobre
\texttt{username}, sobre \texttt{camera\_id} y sobre la matrícula del padrón
(\texttt{authorized\_vehicles.license\_plate}), que impiden duplicados en las
tres entidades de control. Estos índices responden a los patrones de consulta
predominantes del sistema: la recuperación del historial de un vehículo, la
consulta de las sesiones activas para el cálculo de la ocupación, la revisión de
las observaciones asociadas a una visita determinada y el contraste de una
matrícula leída contra el padrón de vehículos habilitados.

\subsection{Modelo de aforo perimetral y gestión de ocupación}

Un aspecto de diseño que delimita el alcance del sistema es el modelo adoptado
para la gestión de ocupación. Frente a los esquemas basados en supervisión de
plazas individuales (\textit{bay monitoring} o \textit{slot guidance}), que
demandan una infraestructura extensiva de cámaras cenitales o sensores de
proximidad en cada cajón, este trabajo adopta el modelo de \textbf{aforo
perimetral}.

Bajo este enfoque, el sistema modela el parqueadero como un recinto delimitado
cuyo estado se gobierna por la ley de balance de flujo continuo:
\begin{equation}
    O(t) = \min\left(C,\, \max\left(0,\, O(t_0) + \sum_{k} I_k(t_0, t) - \sum_{k} E_k(t_0, t)\right)\right)
\end{equation}
donde $O(t)$ representa el número de vehículos activos en el recinto en el
instante $t$, $C$ es la capacidad nominal configurada, e $I_k$ y $E_k$ son los
eventos de ingreso y egreso confirmados por el subsistema de visión artificial
en los puntos perimetrales de acceso $k$. El porcentaje de ocupación relativo se
expresa como $P(t) = (O(t) / C) \times 100$.

Este diseño presenta tres ventajas operativas y técnicas determinantes:
\begin{enumerate}
    \item \textbf{Eficiencia de infraestructura:} requiere cobertura óptica
        exclusivamente en las vías de entrada y salida, reduciendo los costos de
        adquisición, cableado y mantenimiento frente a un esquema de cámaras
        internas.
    \item \textbf{Inmunidad a oclusiones internas:} no depende de líneas de
        visión directas hacia cada cajón, las cuales suelen verse degradadas por
        columnas, techumbres bajas o vehículos de gran porte adyacentes.
    \item \textbf{Trazabilidad vinculada a eventos:} cada unidad computada en el
        aforo corresponde rigurosamente a una sesión abierta con evidencia
        visual (recorte del vehículo y matrícula auditada), evitando discrepancias
        entre sensores dispersos.
\end{enumerate}

\subsection{Especificación de la interfaz de programación}

La API se diseña siguiendo el estilo REST y expone sus servicios bajo un prefijo
versionado, de modo que futuras modificaciones no invaliden a los clientes
existentes. La Tabla \ref{tab:endpoints} especifica los recursos previstos por el
diseño. Nótese que el registro de ingreso y de salida queda reservado a los
dispositivos de captura, mientras que las operaciones de consulta se restringen
a los usuarios autenticados. En esa formulación, los recursos de registro
responden de forma inmediata con el identificador de la sesión creada, sin
esperar a que concluya la inferencia, que se ejecuta de manera asíncrona.

La implementación conserva la separación de responsabilidades, pero desplaza la
ingesta fuera de la interfaz HTTP: el proceso de captura publica la tarea de
inferencia directamente en la cola y el trabajador escribe el resultado en la
base de datos, de modo que no existen recursos de registro de ingreso ni de
salida y el camino de escritura no atraviesa el ciclo de atención de una
petición. El vocabulario de las rutas se unificó en inglés, conforme a la
convención de nomenclatura del código. Los recursos efectivamente expuestos se
enumeran en la Tabla~\ref{tab:api_implementada} del capítulo de implementación.

\begin{longtable}{
        P{0.10}
        P{0.27}
        P{0.38}
        P{0.15}
    }
    \caption{Especificación de los recursos de la API}
    \label{tab:endpoints}\\

    \toprule
    \rowcolor{gray!20}
    \textbf{Método} &
    \textbf{Ruta} &
    \textbf{Descripción} &
    \textbf{Acceso} \\
    \midrule
    \endfirsthead

    \multicolumn{4}{c}{\textit{Continuación de la Tabla
    \ref{tab:endpoints}}} \\
    \toprule
    \rowcolor{gray!20}
    \textbf{Método} &
    \textbf{Ruta} &
    \textbf{Descripción} &
    \textbf{Acceso} \\
    \midrule
    \endhead

    POST & /api/v1/entry &
    Recibe la imagen del vehículo en el acceso, crea la sesión y encola la
    inferencia. &
    Dispositivo de captura \\
    \midrule

    POST & /api/v1/exit &
    Recibe la imagen del vehículo en la salida y solicita el cierre de la
    sesión. &
    Dispositivo de captura \\
    \midrule

    GET & /api/v1/vehiculos &
    Consulta los vehículos registrados con sus atributos visuales y su
    historial. &
    Operador \\
    \midrule

    POST & /api/v1/whitelist &
    Registra o actualiza un vehículo en el padrón del recinto (RF-10). &
    Administrador \\
    \midrule

    GET & /api/v1/stats/active &
    Lista las sesiones actualmente activas dentro del recinto. &
    Operador \\
    \midrule

    GET & /api/v1/stats/history &
    Consulta el historial de sesiones con filtros por fecha. &
    Operador \\
    \midrule

    GET & /api/v1/stats/occupancy &
    Devuelve el nivel de ocupación del recinto. &
    Operador \\
    \midrule

    GET & /api/v1/stats/daily &
    Devuelve los indicadores agregados del día en curso. &
    Operador \\
    \midrule

    WS & /api/v1/ws/occupancy &
    Canal persistente que propaga las variaciones de ocupación en tiempo real. &
    Operador \\*
    \multicolumn{4}{r}{\footnotesize\textit{Fuente: Elaboración propia}} \\*
    \bottomrule
\end{longtable}

La especificación formal de estos recursos se publica de forma automática a
partir de las definiciones de tipo del propio servicio, lo que garantiza que la
documentación de la interfaz permanezca sincronizada con su implementación. La
implementación incorpora además los recursos de gestión de usuarios y de registro
de cámaras, que se documentan en el capítulo de implementación.

\subsection{Procesamiento asíncrono y notificación en tiempo real}

La inferencia de los modelos de visión es intensiva en cómputo y de duración
variable, por lo que no puede ejecutarse dentro del ciclo de atención de una
solicitud HTTP sin comprometer la disponibilidad del servicio. Ambos
requerimientos se resuelven mediante dos mecanismos complementarios.

El primero es la ejecución asíncrona: al recibir una imagen, el servicio crea el
registro de la sesión en estado activo y encola una tarea de inferencia,
respondiendo de inmediato al dispositivo de captura. El trabajador consume la
tarea cuando dispone de capacidad y, al concluir, actualiza la sesión con los
atributos reconocidos. De este modo, una ráfaga de vehículos se traduce en una
cola de tareas pendientes y no en una degradación del tiempo de respuesta de la
API.

El segundo es la notificación en tiempo real. La aplicación web mantiene un
canal persistente con el servicio, a través del cual recibe las variaciones de
ocupación sin necesidad de consultar periódicamente al servidor. Este diseño
reduce el tráfico de red y, sobre todo, elimina el retardo inherente al sondeo:
el panel de monitoreo refleja el cambio en el instante en que el sistema lo
consolida. Ante una interrupción del canal, el cliente restablece la conexión y
reconcilia su estado consultando los recursos de historial y de ocupación.

\subsection{Trazabilidad y control de autorización}

La trazabilidad se garantiza en tres niveles. En el nivel de integridad, las
claves foráneas y las restricciones del esquema impiden la existencia de
sesiones sin vehículo asociado. En el nivel de evidencia, la bitácora de
detecciones conserva cada observación producida por los modelos junto con su
confianza, su latencia de inferencia y el identificador de la versión del modelo
empleado, lo que permite reconstruir con qué versión del sistema se tomó cada
decisión. En el nivel de correlación, cada sesión queda vinculada al vehículo
identificado, a la puerta de acceso y a las detecciones que sustentaron la
lectura de su matrícula.

El contraste de la matrícula reconocida contra un padrón de vehículos
autorizados (RF-10) se materializa mediante la entidad
\textit{vehículo\_autorizado} (\texttt{authorized\_vehicles}), complementando a la
entidad operativa \textit{vehículo} que registra las observaciones de acceso.
Al consolidarse un evento de ingreso con placa reconocida, el sistema verifica
la existencia y vigencia de la matrícula en el padrón institucional, alertando al
operador ante accesos de vehículos no habilitados.

Dado que el trabajador de inferencia puede reintentar la actualización de una
sesión tras una interrupción, la operación de cierre se diseña como idempotente:
la marca temporal de salida y la duración se calculan una única vez, y los
reintentos posteriores sobre una sesión ya cerrada se ignoran. Esta decisión
evita que una falla de comunicación derive en una duración inconsistente o en un
cierre duplicado.

Asimismo, la arquitectura incorpora un mecanismo explícito de tolerancia a fallos
frente a lecturas imprecisas del reconocedor de matrículas. En caso de que la
matrícula leída en un evento de salida no encuentre una coincidencia exacta
activa en la base de datos ---situación esperable ante la tasa de error por
carácter del OCR---, el sistema no bloquea el flujo ni corrompe el estado
transaccional: clasifica el evento como una discrepancia y genera una alerta en la
interfaz operativa de incidentes. De este modo, el operador dispone en pantalla de
los fotogramas de entrada y salida, el recorte de la placa y los registros
candidatos por orden de aproximación temporal, permitiendo realizar el cierre de
sesión de forma manual, controlada y auditada (\texttt{cierre\_operador} y
\texttt{cierre\_motivo}).
